\title{DeepSeekMath: Pushing the Limits of Mathematical Reasoning in Open Language Models}

\begin{abstract}
The inherently intricate and structured characteristics of mathematical reasoning create substantial obstacles for language models. This study presents DeepSeekMath~7B, an extension of DeepSeek-Coder-Base-v1.5 7B that is further pre-trained on 120B math-specific tokens drawn from Common Crawl, in addition to natural language and code resources. Without employing external toolkits or ensemble techniques, DeepSeekMath~7B achieves a strong performance of 51.7\% on the MATH benchmark at the competition level, nearly matching results reported by Gemini-Ultra and GPT-4. Additionally, when utilizing self-consistency sampling with 64 outputs, DeepSeekMath~7B attains 60.9\% accuracy on MATH. The model's robust mathematical reasoning stems from two primary elements: firstly, we tap into the vast potential of publicly accessible web data by implementing a rigorously designed data filtering pipeline; secondly, we propose Group Relative Policy Optimization (GRPO)—an adaptation of Proximal Policy Optimization (PPO)—which augments mathematical reasoning skills and enhances PPO's memory efficiency in parallel.
\end{abstract}

\section{Introduction}

Large language models (LLMs) have transformed artificial intelligence methodologies for mathematical reasoning, leading to notable progress in benchmarks for quantitative reasoning \citep{MATH} as well as geometry-based tasks \citep{trinh2024solving}. These systems have also demonstrated value in supporting humans when tackling intricate mathematical challenges \citep{tao}. Despite these achievements, state-of-the-art architectures like GPT-4 \citep{gpt4} and Gemini-Ultra \citep{gemini} remain unavailable for public use, while extant open-source models exhibit substantial performance gaps in comparison.

This paper presents DeepSeekMath, a language model tailored for mathematical reasoning that surpasses the mathematical performance of existing open-source systems and nearly matches GPT-4 on standard academic assessments. Central to this achievement is the construction of the DeepSeekMath Corpus, an extensive, high-quality pre-training dataset containing 120B math-specific tokens. The corpus is sourced from Common Crawl (CC) with the assistance of a fastText-based classifier \citep{joulin2016fasttext}. The initial version of this classifier is trained using OpenWebMath \citep{openwebmath} samples as positive cases, while negatives are drawn from a broad array of web documents. The classifier then identifies new candidate math instances from CC, which are curated and verified by human annotators. Incorporating these refined samples, the classifier is iteratively improved on an enriched training set. Empirical evaluation shows that this large-scale dataset leads to strong downstream performance, with DeepSeekMath-Base 7B attaining scores of 64.2\% on GSM8K \citep{gsm8k} and 36.2\% on the advanced-level MATH dataset \citep{MATH}, outperforming Minerva 540B \citep{minerva}. Furthermore, as the DeepSeekMath Corpus includes multilingual content, we observe enhanced results on Chinese mathematics benchmarks \citep{wei2023cmath,agieval}. We propose that our methodology in curating mathematical corpora lays a foundation for further studies and that substantial advances remain possible.

DeepSeekMath-Base is built upon DeepSeek-Coder-Base-v1.5 7B \citep{deepseek-coder}, based on our observation that initializing from a code-oriented pretrained model yields superior results than beginning with a standard language model. Additionally, our findings show that mathematical training benefits the model’s performance on the MMLU \citep{mmlu} and BBH \citep{bbh} benchmarks, suggesting that such training bolsters both mathematical proficiency and broader reasoning skills.

Following the pre-training phase, we conduct mathematical instruction tuning on DeepSeekMath-Base utilizing datasets comprising chain-of-thought \citep{cot}, program-of-thought \citep{pot,pal}, and tool-integrated reasoning \citep{tora} examples. The tuned model, DeepSeekMath-Instruct 7B, surpasses all other 7B models and demonstrates performance on par with open-source instruction-tuned models at the 70B scale.

In addition, we propose Group Relative Policy Optimization (GRPO), an alternative reinforcement learning (RL) method based on Proximal Policy Optimization (PPO) \citep{schulman2017proximal}. Unlike PPO, GRPO eliminates the need for a critic model, leveraging group scores as a baseline and thereby markedly decreasing computational demands during training. Utilizing only a portion of English instruction tuning data, GRPO delivers notable gains over the competitive DeepSeekMath-Instruct system, improving results on both in-domain benchmarks (GSM8K: 82.9\% $\rightarrow$ 88.2\%, MATH: 46.8\% $\rightarrow$ 51.7\%) and out-of-domain mathematical evaluation (e.g., CMATH: 84.6\% $\rightarrow$ 88.8\%) in the RL stage. We further present a cohesive framework for interpreting various methods, including Rejection Sampling Fine-Tuning (RFT) \citep{yuan2023scaling}, Direct Preference Optimization (DPO) \citep{dpo}, PPO, and GRPO. With this framework, we reveal that all these approaches can be viewed as instances of direct or streamlined RL algorithms. To comprehensively examine the foundational components of this approach, we carry out rigorous empirical studies addressing aspects such as online vs. offline learning, outcome-based vs. process-oriented supervision, and single-step vs. iterative RL. Finally, we discuss the mechanisms by which RL enhances instruction-tuned models and identify prospective avenues for advancing RL effectiveness within this unified conceptual landscape.

\subsection{Contributions}
We present a scalable approach to math pre-training, and further conduct a detailed investigation and assessment of reinforcement learning techniques.

\textbf{Math Pre-Training at Scale}

\begin{itemize}[topsep=0pt]
    \item Our study demonstrates that the Common Crawl dataset, which is publicly available, offers significant resources pertinent to mathematics. Through the use of a carefully engineered data selection pipeline, we assembled the DeepSeekMath Corpus, a high-quality collection comprising 120B tokens from mathematically relevant web pages. This corpus is nearly seven times larger than the math web dataset utilized by Minerva \citep{minerva} and nine times bigger than the recently introduced OpenWebMath corpus \citep{openwebmath}.

\item DeepSeekMath-Base 7B, our pre-trained base model, attains performance similar to Minerva 540B \citep{minerva}, suggesting that parameter count alone does not solely determine mathematical reasoning proficiency. High-quality pre-training data enables smaller models to perform at a competitive level.

\item We present the outcomes of our experiments involving math training. Exposing models to code training before commencing math training enhances their performance in solving mathematical tasks, irrespective of whether external tools are used. This provides some evidence toward resolving the enduring inquiry: \textit{does code training bolster reasoning skills?} Our results suggest that, in the context of mathematical reasoning, it does.

\item While the use of arXiv papers as training data is prevalent—particularly within mathematics-focused literature—it does not yield any significant performance gains on the mathematical benchmarks considered in this study.
\end{itemize}

\textbf{Exploration and Analysis of Reinforcement Learning}

\begin{itemize}[topsep=0pt]
    \item We present Group Relative Policy Optimization (GRPO), a reinforcement learning method characterized by both high efficiency and effectiveness. Unlike traditional approaches, GRPO dispenses with the critic network, opting instead to derive the baseline from aggregated group scores. This modification leads to a considerable decrease in computational and training resource demands compared to Proximal Policy Optimization (PPO).
    \item Through empirical evaluation, we show that GRPO provides substantial improvements for our instruction-tuned model DeepSeekMath-Instruct using exclusively instruction-tuning datasets. Notably, the reinforcement learning phase also brings gains in the generalization capabilities of the model to domains not represented in the training set.
    \item We establish a cohesive conceptual framework for interpreting a variety of methods—specifically RFT, DPO, PPO, and GRPO. To further analyze this unified perspective, we perform comprehensive experiments, examining aspects such as online versus offline learning, the comparative impact of outcome and process supervision, as well as distinctions between single-step and iterative forms of reinforcement learning, among other factors.
    \item Leveraging this integrative framework, we delve into the underlying factors that contribute to the success of reinforcement learning, and we outline several promising pathways for advancing more robust reinforcement learning approaches for large language models (LLMs).
\end{itemize}

\subsection{Summary of Evaluations and Metrics}

\begin{itemize}[topsep=0pt]
    \item \textbf{English and Chinese Mathematical Reasoning}: We carry out thorough evaluations of our models across both English and Chinese datasets that span mathematical tasks from elementary to collegiate complexity. The English evaluation suite consists of GSM8K \citep{gsm8k}, MATH \citep{MATH}, SAT \citep{llemma}, OCW Courses \citep{minerva}, and MMLU-STEM \citep{mmlu}. For Chinese, we utilize benchmarks such as MGSM-zh \citep{mgsm}, CMATH \citep{wei2023cmath}, Gaokao-MathCloze \citep{agieval}, and Gaokao-MathQA \citep{agieval}. The assessments focus both on the capacity of models to independently compose complete textual solutions and on their proficiency in leveraging Python for problem-solving.

When evaluated on English benchmarks, DeepSeekMath-Base demonstrates performance comparable to the proprietary Minerva 540B model \citep{minerva}, and consistently outperforms all open-source base models—including Mistral 7B \citep{mistral} and Llemma-34B \citep{llemma}—regardless of their math pre-training status, frequently by a notable margin. Particularly on Chinese benchmarks, DeepSeekMath-Base achieves even greater superiority. This is plausibly due to our departure from earlier approaches \citep{minerva,llemma} that relied exclusively on English mathematical pre-training data, as we incorporate substantial high-quality data from other languages as well. Through the application of mathematical instruction tuning and reinforcement learning, the resultant DeepSeekMath-Instruct and DeepSeekMath-RL models exhibit robust performance, notably surpassing the 50\% accuracy threshold for the first time on the competition-level MATH dataset among open-source efforts.

\item \textbf{Formal Mathematics}: We assess DeepSeekMath-Base's abilities in informal-to-formal theorem translation as outlined in \citep{dsp_proof}, utilizing the miniF2F dataset \citep{minif2f} with Isabelle \citep{isabelle} serving as the proof assistant. Notably, DeepSeekMath-Base achieves robust performance in few-shot autoformalization scenarios.
\item \textbf{Natural Language Understanding, Reasoning, and Code}: For a holistic assessment of the models' proficiency in general comprehension, complex reasoning, and programming, we administer evaluations on the Massive Multitask Language Understanding (MMLU) benchmark \citep{mmlu}, which contains 57 multiple-choice tasks from various domains. We also test on BIG-Bench Hard (BBH) \citep{bbh}, a suite of 23 tasks largely centered on multi-step reasoning, alongside HumanEval \citep{codex} and MBPP \citep{mbpp}, two established benchmarks for code generation. Pre-training on mathematical data enhances both language comprehension and reasoning skills.

\section{Math Pre-Training}

\subsection{Data Collection and Decontamination} This section details the methodology employed to build the DeepSeekMath Corpus utilizing data sourced from Common Crawl. As illustrated in Figure \ref{fig:data_collect}, we describe an iterative workflow designed to methodically compile an extensive mathematical dataset from Common Crawl, beginning with an initial seed set (for example, a limited yet high-quality repository of mathematics-related data). It should be emphasized that this methodology can likewise be extended to additional domains, including programming.

To begin, we select OpenWebMath \citep{openwebmath}, a curated repository of high-quality mathematics-oriented web documents, as our primary seed dataset. Using this data, we train a fastText model \citep{joulin2016fasttext} to retrieve additional mathematical web pages resembling those in OpenWebMath. For this purpose, we randomly sample 500,000 entries from the seed corpus to serve as positive examples, while another 500,000 pages drawn from Common Crawl act as negative samples. Model training utilizes an open-source implementation\footnote{\url{https://fasttext.cc}}, with hyperparameters set as follows: 256-dimensional vectors, a learning rate of 0.1, a word n-gram window of up to 3, a minimum word frequency threshold of 3, and 3 epochs. To pare down the scope of Common Crawl, we apply both URL-based deduplication and near-duplicate detection, thereby obtaining a collection of 40B unique HTML web pages. Using the trained fastText model, we identify mathematical content within this deduplicated corpus. To exclude low-quality pages, we sort candidates by their fastText model scores and retain only those with the highest rankings. We evaluate the scale of preserved content by performing pre-training trials at the 40B, 80B, 120B, and 160B token levels. For the initial round, we opt to keep the leading 40B tokens.

Following the initial round of data acquisition, a significant number of mathematical web pages remain excluded, largely due to the limited diversity among positive samples used to train the fastText model. To address this, we broaden the scope of seed material by sourcing additional mathematical websites, aiming to enhance the effectiveness of the fastText model. Specifically, the complete Common Crawl dataset is partitioned according to discrete domains, with each domain encompassing all pages under the same base URL. For every domain, we determine what fraction of its web pages were included during the first round of collection. Domains where more than 10\% of pages were captured are labeled as math-related (for example, \url{mathoverflow.net}). 

Next, we conduct manual annotation of URLs within these targeted domains to identify those specifically associated with mathematical content (e.g., \url{mathoverflow.net/questions}). Web pages that are linked to these URLs but have not yet been collected are subsequently incorporated into the seed corpus. This strategy yields a richer pool of positive training examples, facilitating the development of an improved fastText model with enhanced recall capability for mathematical content in subsequent collection cycles. Across four iterations of this process, we obtain a total of 35.5M mathematical web pages, amassing 120B tokens. By the fourth iteration, analysis reveals that about 98\% of the mathematical data had already been harvested during the third round, prompting us to terminate the data collection effort.

In order to prevent benchmark contamination, we adopt the filtering approach described in \cite{deepseek-coder}, removing any web pages that include questions or answers from established English mathematical benchmarks like GSM8K \citep{gsm8k} and MATH \citep{MATH}, as well as Chinese datasets such as CMATH \citep{wei2023cmath} and AGIEval \citep{agieval}. Specifically, our filtering process excludes from the math training data any text span that exactly matches a 10-gram sequence from the benchmark datasets. In cases where benchmark texts consist of fewer than 10 grams but include at least 3-gram sequences, those shorter spans are also matched exactly and the corresponding web pages are filtered out if they are found.

\subsection{Validating the Quality of the DeepSeekMath~Corpus}

We run pre-training experiments to investigate how the DeepSeekMath~Corpus is compared with the recently released math-training corpora:

\begin{itemize}[topsep=0pt]
    \item \textbf{MathPile} \citep{mathpile}: This collection brings together 8.9B tokens from diverse sources, such as textbooks, Wikipedia, ProofWiki, CommonCrawl, StackExchange, and notably arXiv, from which over 85\% of the data is derived;
    \item \textbf{OpenWebMath} \citep{openwebmath}: Curated from CommonCrawl by extracting materials with mathematical content, this dataset comprises 13.6B tokens;
    \item \textbf{Proof-Pile-2} \citep{llemma}: An extensive mathematical dataset made up of OpenWebMath, AlgebraicStack (containing 10.3B tokens of math-centric code), and 28.0B tokens from arXiv publications. In our experiments using Proof-Pile-2, we adhere to the arXiv:Web:Code distribution of 2:4:1, as established by \cite{llemma}.
\end{itemize}

\subsubsection{Training Setting}
\label{sec:quality-policy}
We fine-tune a general-purpose language model with 1.3B parameters on mathematical data, where the model architecture mirrors that of the DeepSeek LLMs \citep{deepseek-llm}, referred to in this work as DeepSeek-LLM 1.3B. For each distinct mathematical dataset, separate models are trained for a total of 150B tokens. All model training is performed with the efficient, lightweight HAI-LLM \citep{haillm} training framework. In alignment with DeepSeek LLMs' training protocol, we employ the AdamW optimizer \citep{adamW} with $\beta_1=0.9$, $\beta_2=0.95$, and a weight decay parameter of $\mathrm{weight\_decay}=0.1$. The learning rate follows a multi-stage schedule: it linearly increases to its maximum after 2,000 warmup steps, declines to 31.6\% of its maximum value after 80\% of the training duration, and then drops further to 10.0\% at the 90\% training mark. The maximum learning rate is set at 5.3e-4. Each training step processes a batch of 4 million tokens and utilizes a 4K token context window.

\subsubsection{Evaluation Results}

\textbf{The DeepSeekMath~Corpus is of high quality, covers multilingual mathematical content, and is the largest in size.}

\begin{itemize}[topsep=0pt]
    \item \textbf{High-quality}:
    We assess the impact on downstream tasks across eight mathematical benchmarks using few-shot chain-of-thought prompting as in \cite{cot}. Table \ref{tab:corpora_comparison} clearly illustrates the superior performance of models trained with the DeepSeekMath~Corpus. Furthermore, Figure \ref{fig:corpora_comparisons} reveals that the DeepSeekMath Corpus outperforms Proof-Pile-2 even after exposure to 50B tokens (corresponding to a complete epoch over Proof-Pile-2), signifying that DeepSeekMath offers higher overall data quality.
    \item \textbf{Multilingual}:
    DeepSeekMath~Corpus contains data from a variety of languages, with English and Chinese being the most prevalent. According to Table \ref{tab:corpora_comparison}, models trained using DeepSeekMath~Corpus show enhanced mathematical reasoning in both English and Chinese. On the other hand, prevailing mathematical datasets—primarily English-focused—result in modest, or even detrimental, effects on Chinese mathematical reasoning tasks.
    \item \textbf{Large-scale}:
    The volume of DeepSeekMath~Corpus significantly exceeds that of previously available mathematical datasets. As illustrated in Figure \ref{fig:corpora_comparisons}, models like DeepSeek-LLM 1.3B, trained with this corpus, experience sharper learning trajectories and sustain their performance gains for a longer period. In comparison, the baseline datasets are considerably smaller and have been cycled through several times during training, leading to a rapid stagnation in model progress.
\end{itemize}

\subsection{Training and Evaluating DeepSeekMath-Base 7B}

This section presents DeepSeekMath-Base 7B, a foundational model characterized by robust mathematical reasoning capabilities. The model builds upon DeepSeek-Coder-Base-v1.5 7B \citep{deepseek-coder}, undergoing training for a total of 500B tokens. Data utilized during training was composed of 56\% DeepSeekMath Corpus, 4\% AlgebraicStack, 10\% arXiv, 20\% code from Github, and 10\% natural language data drawn from Common Crawl in both English and Chinese. Training configurations predominantly follow the approach detailed in Section \ref{sec:quality-policy}, with the exceptions of employing a maximum learning rate of 4.2e-4 and a batch size of 10M tokens.

This study provides an in-depth evaluation of DeepSeekMath-Base 7B's mathematical proficiency, examining its performance in generating self-sufficient mathematical solutions without tool assistance, its effectiveness in addressing mathematical problems with tool support, and its competency in formal theorem proving tasks. In addition to mathematics-specific analysis, we further outline a broader characterization of the base model by examining its capabilities in natural language comprehension, reasoning, and programming.

\paragraph{Mathematical Problem Solving with Step-by-Step Reasoning}

We assess DeepSeekMath-Base’s capability to address mathematical questions through few-shot chain-of-thought prompting \citep{cot}, spanning eight distinct benchmarks in both English and Chinese. The selected benchmarks address a broad spectrum of mathematical domains, incorporating tasks focused on quantitative reasoning (such as GSM8K \citep{gsm8k}, MATH \citep{MATH}, and CMATH \citep{wei2023cmath}) as well as multiple-choice formats (including MMLU-STEM \citep{mmlu} and Gaokao-MathQA \citep{agieval}). Together, these benchmarks test understanding from basic to advanced undergraduate mathematics.

Table \ref{tab:base_math_cot} demonstrates that among open-source base models—encompassing prominent models such as Mistral 7B \citep{mistral} and Llemma 34B \citep{llemma}, the latter of which is enhanced by math-specific training on the Proof-Pile-2 corpus \citep{llemma}—DeepSeekMath-Base 7B consistently attains top performance across all eight benchmarks. In particular, for the challenging competition-grade MATH dataset, DeepSeekMath-Base exceeds the accuracy of all current open-source base models by a margin exceeding 10\% in absolute terms. Furthermore, it also outperforms Minerva 540B \citep{minerva}, a substantially larger (77-fold) proprietary base model that is developed atop PaLM \citep{palm} and further fine-tuned on specialized mathematical literature.

\paragraph{Mathematical Problem Solving with Tool Use} We assess the capacity for mathematical reasoning with the assistance of external tools on the GSM8K and MATH benchmarks, employing few-shot program-of-thought prompting as described in \citep{pot,pal}. In this setup, models address each question by generating a Python script, making use of libraries such as \textit{math} and \textit{sympy} to handle complex calculations. The output generated from running these scripts is treated as the predicted answer. According to Table \ref{tab:base_math_pot_proof}, DeepSeekMath-Base 7B demonstrates superior performance over the previous best, Llemma 34B.

\paragraph{Formal Mathematics}

Automating the process of formal proof verification offers significant advantages for ensuring both the correctness and dependability of mathematical arguments, as well as improving productivity—a topic that has seen growing interest in recent years. In this work, we assess the capabilities of DeepSeekMath-Base 7B on the informal-to-formal proving task from \citep{dsp_proof}. This task involves constructing a formal proof given an informal problem statement, its corresponding formal version, and an informal proof. Our evaluation employs the miniF2F benchmark \citep{minif2f}, which focuses on formalizations of Olympiad-caliber mathematics. For each problem, we use few-shot prompting to produce an Isabelle formal proof. Adhering to the methodology of \cite{dsp_proof}, the model is tasked with creating proof outlines, after which the Sledgehammer automated prover \citep{sledgehammer} is employed to address incomplete components. According to Table \ref{tab:base_math_pot_proof}, DeepSeekMath-Base 7B achieves notable results in the automatic formalization of proofs.

\paragraph{Natural Language Understanding, Reasoning, and Code}

We assess natural language understanding using MMLU \citep{mmlu}, evaluate reasoning with BBH \citep{bbh}, and test coding proficiency on HumanEval \citep{codex} and MBPP \citep{mbpp}. As reported in Table \ref{tab:understanding_reasoning_code}, DeepSeekMath-Base 7B demonstrates marked improvements over its predecessor, DeepSeek-Coder-Base-v1.5 \citep{deepseek-coder}, in both MMLU and BBH, highlighting the beneficial effects of mathematical training on language comprehension and reasoning skills. Moreover, through the integration of code tokens during continual training, DeepSeekMath-Base 7B successfully retains the coding performance seen in DeepSeek-Coder-Base-v1.5 on the coding benchmarks. In summary, DeepSeekMath-Base 7B achieves substantially better results than the general-purpose Mistral 7B \citep{mistral} across these reasoning and coding evaluations.

\section{Supervised Fine-Tuning}

\subsection{SFT Data Curation}

We develop a dataset for mathematical instruction-tuning that spans both English and Chinese, encompassing a broad spectrum of mathematical domains and degrees of difficulty. Each problem in the dataset is matched with corresponding solutions in chain-of-thought (CoT) \citep{cot}, program-of-thought (PoT) \citep{pot,pal}, and tool-integrated reasoning \citep{tora} formats. In total, the dataset comprises 776K training instances.

\begin{itemize}[topsep=0pt]
    \item \textbf{English mathematical datasets}: We provide annotated solutions that utilize tool integration for problems from GSM8K and MATH, in addition to incorporating a selected subset of MathInstruct \citep{MathInstruct} and the training portion of Lila-OOD \citep{lila}, in which questions are resolved using either CoT or PoT techniques. Our assembled English dataset encompasses a broad range of mathematical disciplines, including areas like algebra, probability, number theory, calculus, and geometry.
    \item \textbf{Chinese mathematical datasets}: We assemble a corpus of K-12 mathematics questions in Chinese, covering 76 distinct sub-domains such as linear equations. Each problem is accompanied by detailed solutions presented in both CoT and tool-integrated reasoning formats.
\end{itemize}

\subsection{Training and Evaluating DeepSeekMath-Instruct 7B}

This section presents DeepSeekMath-Instruct 7B, which is developed through mathematical instruction tuning derived from DeepSeekMath-Base. To prepare training data, examples are randomly joined together until the total length reaches up to 4K tokens. The model is then trained over 500 steps, utilizing a batch size of 256 and maintaining a fixed learning rate of 5e-5 throughout.

We evaluate models' mathematical performance both without and with tool use, on 4 quantitative reasoning benchmarks in English and Chinese. We benchmark our model against the leading models of the time:

\begin{itemize}[topsep=0pt]
    \item \textbf{Closed-source models} comprise: (1) the GPT series, particularly GPT-4 \citep{gpt4} and GPT-4 Code Interpreter~\footnote{\url{https://openai.com/blog/chatgpt-plugins\#\#code-interpreter}}, which demonstrate leading capabilities; (2) Gemini Ultra and Pro \citep{gemini}; (3) Inflection-2 \citep{inflection-2}; (4) Grok-1~\footnote{\url{https://x.ai/model-card}}; as well as newer offerings from Chinese technology companies, such as (5) Baichuan-3~\footnote{\url{https://www.baichuan-ai.com}}, and (6) the most recent GLM-4~\footnote{\url{https://open.bigmodel.cn/dev/api\#glm-4}}.

    from the GLM family \citep{glm}.

These models serve broad purposes and typically have been subjected to extensive alignment processes.

\item \textbf{Open-source models} feature both general-purpose options such as (1) DeepSeek-LLM-Chat 67B \citep{deepseek-llm}, (2) Qwen 72B \citep{qwen}, (3) SeaLLM-v2 7B \citep{seallm}, and (4) ChatGLM3 6B \citep{chatglm3}, alongside models specifically optimized for mathematical performance, including: (5) InternLM2-Math 20B~\footnote{\url{https://github.com/InternLM/InternLM-Math}}, a variant built atop InternLM2 with additional math-specific pre-training and subsequent instruction tuning; (6) Math-Shepherd-Mistral 7B, which incorporates PPO-based training \citep{schulman2017proximal} and leverages a process-supervised reward model applied to the Mistral 7B architecture \citep{mistral}; (7) the WizardMath family \citep{wizardmath}, focused on enhancing the mathematical reasoning capabilities of Mistral 7B and Llama-2 70B \citep{llama2} via evolve-instruct methods—i.e., instruction tuning enriched with AI-generated instructions—and PPO training primarily on examples from GSM8K and MATH; (8) MetaMath 70B \citep{metamath}, consisting of Llama-2 70B further trained on expanded datasets derived from GSM8K and MATH; (9) ToRA 34B \cite{tora}, which is CodeLlama 34B with fine-tuning targeted at integrating tool-augmented mathematical problem solving; and (10) MAmmoTH 70B \citep{MathInstruct}, representing Llama-2 70B adapted through MathInstruct-based instruction tuning.

Referencing Table \ref{tab:sft_rl_math}, when evaluated without permitting tool usage, DeepSeekMath-Instruct 7B exhibits notable capability in sequential reasoning tasks. Specifically, on the competition-level MATH dataset, our model exceeds the performance of all open-source counterparts and outperforms most proprietary systems—such as Inflection-2 and Gemini Pro—by a margin of at least 9\% in absolute terms. This advantage persists even compared to models with considerably larger parameter counts (e.g., Qwen 72B) or those that have undergone specialized math-centric reinforcement learning (such as WizardMath-v1.1 7B). Although DeepSeekMath-Instruct achieves results comparable to leading Chinese proprietary models like GLM-4 and Baichuan-3, its accuracy still lags behind top-tier models, including GPT-4 and Gemini Ultra.

In an evaluation scenario permitting models to combine natural language reasoning with program-driven tool usage for solving problems, DeepSeekMath-Instruct 7B attains nearly 60\% accuracy on the MATH dataset, outperforming all previously available open-source models. Across additional benchmarks, our model matches the performance of DeepSeek-LLM-Chat 67B, the previous state-of-the-art despite its model size being tenfold larger.

\subsection{Group Relative Policy Optimization} Following the Supervised Fine-Tuning (SFT) phase, reinforcement learning (RL) methods have demonstrated significant improvements in the mathematical reasoning capacities of LLMs \citep{wang2023math,wizardmath}. Here, we present our own approach, Group Relative Policy Optimization (GRPO), an RL algorithm designed to offer both efficiency and effectiveness.

\subsubsection{From PPO to GRPO} Proximal Policy Optimization (PPO) \citep{schulman2017proximal} is an actor-critic RL algorithm that is widely used in the RL fine-tuning stage of LLMs \citep{ouyang2022training}. In particular, it optimizes LLMs by maximizing the following surrogate objective:

\begin{equation}
\footnotesize
    \mathcal{J}_{PPO}(\theta) = \mathbb{E}{[q \sim P(Q), o \sim \pi_{\theta_{old}}(O|q)]} \frac{1}{|o|} \sum_{t=1}^{|o|} \min \left[ \frac{\pi_\theta(o_{t} | q, o_{<t})}{\pi_{\theta_{old}}(o_{t} | q, o_{<t})} A_{t}, \text{clip} \left( \frac{\pi_\theta(o_{t} | q, o_{<t})}{\pi_{\theta_{old}}(o_{t} | q, o_{<t})}, 1 - \varepsilon, 1 + \varepsilon \right)  A_{t} \right] ,
\end{equation}

The above equation presents the objective function for Proximal Policy Optimization (PPO). This loss, $\mathcal{J}_{PPO}(\theta)$, is calculated as the expectation over queries $q$ sampled from $P(Q)$ and outputs $o$ drawn according to the behavior policy $\pi_{\theta_{old}}(O|q)$. For each output sequence $o$, the loss is averaged across its time steps. At each time step $t$, the minimum is taken between the ratio of the current and old policies, scaled by the advantage $A_t$, and this same ratio clipped within $[1-\varepsilon, 1+\varepsilon]$ before scaling by $A_t$. This approach serves to mitigate large policy updates during training.

Here, $\pi_{\theta}$ denotes the current policy, while $\pi_{\theta_{old}}$ refers to the previous policy iteration. The variables $q$ and $o$ correspond to questions drawn from the question dataset and outputs generated by $\pi_{\theta_{old}}$, respectively. The hyper-parameter $\varepsilon$ is employed in PPO as part of the clipping mechanism designed to promote stable learning dynamics. The term $A_t$ represents the advantage estimate, which is obtained using Generalized Advantage Estimation (GAE) \citep{gae}, relying on future rewards $\{r_{\ge t}\}$ in conjunction with a value function $V_{\psi}$ that is learned during training. Consequently, PPO requires the simultaneous training of a value function along with the policy. To prevent excessive optimization with respect to the reward model, a common strategy is to introduce a per-token KL divergence penalty, computed with respect to a reference model and integrated into the token-level reward signal \citep{ouyang2022training}, as follows:

\begin{equation}
    r_{t} = r_\varphi(q, o_{\le t}) - \beta \log\frac{\pi_{\theta}(o_{t}|q, o_{<t})}{\pi_{ref}(o_{t}|q, o_{<t})},
\label{eq:PPO-reward}
\end{equation}

Here, $r_\varphi$ denotes the reward model, $\pi_{ref}$ refers to the reference model—commonly the original SFT model—and $\beta$ represents the weight assigned to the KL divergence penalty.

As the value function employed in PPO is typically another model of comparable size as the policy model, it brings a substantial memory and computational burden. Additionally, during RL training, the value function is treated as a baseline in the calculation of the advantage for variance reduction. While in the LLM context, usually only the last token is assigned a reward score by the reward model, which may complicate the training of a value function that is accurate at each token. To address this, as shown in Figure \ref{fig:grpo}, we propose Group Relative Policy Optimization (GRPO), which obviates the need for additional value function approximation as in PPO, and instead uses the average reward of multiple sampled outputs, produced in response to the same question, as the baseline. More specifically, for each question $q$, GRPO samples a group of outputs $\{o_1, o_2, \cdots, o_G\}$  from the old policy  $\pi_{\theta_{old}}$  and then optimizes the policy model by maximizing the following objective:

\begin{equation}
\footnotesize
\begin{split}
    \mathcal{J}_{GRPO}(\theta) &= \mathbb{E}{[q \sim P(Q), \{o_i\}_{i=1}^G \sim \pi_{\theta_{old}}(O|q)]}  \\
    & \frac{1}{G}\sum_{i=1}^G\frac{1}{|o_i|} \sum_{t=1}^{|o_i|} \left\{ \min \left[ \frac{\pi_\theta(o_{i,t} | q, o_{i,<t})}{\pi_{\theta_{old}}(o_{i,t} | q, o_{i,<t})} \hat{A}_{i,t}, \text{clip} \left( \frac{\pi_\theta(o_{i,t} | q, o_{i,<t})}{\pi_{\theta_{old}}(o_{i,t} | q, o_{i,<t})}, 1 - \varepsilon, 1 + \varepsilon \right)  \hat{A}_{i,t} \right] - \beta \mathbb{D}_{KL}\left[\pi_{\theta} || \pi_{ref}\right]\right\} ,
\end{split}
\label{eq:GRPO-obj}
\end{equation}

where $\varepsilon$ and $\beta$ are hyper-parameters, and $\hat{A}_{i,t}$  is the advantage calculated based on relative rewards of the outputs inside each group only, which will be detailed in the following subsections. The group relative way that GRPO leverages to calculate the advantages, aligns well with the comparative nature of rewards models,  as reward models are typically trained on datasets of comparisons between outputs on the same question. Also note that, instead of adding KL penalty in the reward, GRPO regularizes by directly adding the KL divergence between the trained policy and the reference policy to the loss, avoiding complicating the calculation of  $\hat{A}_{i,t}$. And different from the KL penalty term used in (\ref{eq:PPO-reward}), we estimate the KL divergence with the following unbiased estimator \citep{kl_approx}:

\begin{equation}
\small
    \mathbb{D}_{KL}\left[\pi_{\theta} || \pi_{ref}\right] = \frac{\pi_{ref}(o_{i,t}|q,o_{i,<t})}{\pi_{\theta}(o_{i,t}|q,o_{i,<t})} - \log\left(\frac{\pi_{ref}(o_{i,t}|q,o_{i,<t})}{\pi_{\theta}(o_{i,t}|q,o_{i,<t})}\right) - 1,
\end{equation}

which is guaranteed to be positive.

\begin{algorithm}[t]
  \small
  \caption{Iterative Group Relative Policy Optimization}
  \textbf{Input} initial policy network $\pi_{\theta_{\text{init}}}$; reward evaluators $r_\varphi$; set of task prompts $\mathcal{D}$;
  hyperparameters $\varepsilon$, $\beta$, $\mu$
  \begin{algorithmic}[1]
    \State Initialize policy network: $\pi_\theta \leftarrow \pi_{\theta_{\text{init}}}$
    \For{iteration = 1, \dots, I}
       \State Set reference model: $\pi_{ref} \leftarrow \pi_{\theta}$
      \For{step = 1, \dots, M}
      \State Draw minibatch $\mathcal{D}_b$ from $\mathcal{D}$
      \State Assign previous policy: $\pi_{\theta_{old}} \leftarrow \pi_{\theta}$
      \State For every prompt $q \in \mathcal{D}_b$, generate $G$ outputs $\{o_i\}_{i=1}^G \sim \pi_{\theta_{old}} (\cdot \mid q) $
      \State For each $o_i$, compute reward $r_i$ using $r_{\varphi}$
      \State Calculate group-relative advantage $\hat{A}_{i,t}$ at the $t$-th token in $o_i$
      \For{GRPO iteration = 1, \dots, $\mu$}
        \State Improve policy $\pi_{\theta}$ by optimizing the GRPO objective (see Equation \ref{eq:GC-GRPO})
      \EndFor
    \EndFor
    \State Continuously retrain $r_\varphi$ using data replay techniques
    \EndFor
  \end{algorithmic}
  \textbf{Output} $\pi_\theta$
  \label{alg:iter-grpo}
\end{algorithm}

\subsubsection{Outcome Supervision RL with GRPO} More specifically, given a question $q$, a set of candidate responses $\{o_1, o_2, \cdots, o_G\}$ are generated using the previous iteration's policy $\pi_{\theta_{old}}$. Each candidate output is evaluated by a reward model, resulting in $G$ associated reward values $\mathbf{r}=\{r_1, r_2, \cdots, r_G\}$. These reward scores are standardized by computing their z-scores, i.e., subtracting the mean of $\mathbf{r}$ and dividing by their standard deviation. Under outcome supervision, each output $o_i$ receives its own normalized reward at its conclusion. For all tokens within a given output, the advantage $\hat{A}_{i, t}$ is uniformly set to the normalized reward: $\hat{A}_{i, t} = \widetilde{r}_i = \frac{r_i- {\rm mean}(\mathbf{r})}{{\rm std}(\mathbf{r})}$. The policy model is then updated by maximizing the objective described in equation (\ref{eq:GRPO-obj}).

\subsubsection{Process Supervision RL with GRPO} 

Outcome supervision only provides feedback in the form of a reward after an entire output sequence is completed, which can be inadequate and inefficient for training policies on intricate mathematical problems. Building on the approach in \cite{wang2023math}, we investigate process supervision, wherein feedback is offered at the conclusion of each individual reasoning step. Specifically, given a question $q$ and $G$ sampled outputs $\{o_1, o_2, \cdots, o_G\}$, a process reward model evaluates each reasoning step, assigning corresponding rewards as follows: $\mathbf{R} = \{ \{r_1^{index(1)},\cdots,r_1^{index(K_1)}\}, \cdots,  \{r_G^{index(1)},\cdots,r_G^{index(K_G)}\} \}$. Here, $index(j)$ denotes the token index where the $j$-th reasoning step ends, and $K_i$ represents the total number of reasoning steps for the $i$-th output. To standardize the range of these rewards, normalization is performed by subtracting the mean and dividing by the standard deviation: $\widetilde{r}_i^{index(j)} = \frac{r_i^{index(j)} - {\rm mean(\mathbf{R})}}{{\rm std(\mathbf{R})}}$.

Following this, process supervision determines the advantage for each token by accumulating the normalized rewards from that token forward, defined as $\hat{A}_{i, t} = \sum_{index(j) \ge t} \widetilde{r}_i^{index(j)}$. The policy is then refined by maximizing the objective specified in equation (\ref{eq:GRPO-obj}).

\subsubsection{Iterative RL with GRPO} During reinforcement learning, the previously trained reward model may become inadequate for guiding the evolving policy model. To address this, we investigate an iterative RL framework utilizing GRPO. As illustrated in Algorithm \ref{alg:iter-grpo}, the iterative GRPO approach involves constructing fresh training datasets for the reward model using samples produced by the policy model, while simultaneously updating the existing reward model with a replay strategy that retains 10\% of past data. Subsequently, the policy model is set as the reference model, after which the policy model is further trained leveraging the newly updated reward model.

\subsection{Training and Evaluating DeepSeekMath-RL}

Our reinforcement learning experiments are performed using the DeepSeekMath-Instruct 7B model. The RL training dataset comprises approximately 144,000 chain-of-thought questions sourced from the SFT dataset, specifically those aligned with GSM8K and MATH, while omitting other SFT entries. This selective approach aims to clarify the effects of RL on benchmarks that do not have overlapping data in the RL training stage. The reward model training data is prepared in accordance with the methodology outlined in \citep{wang2023math}. The reward model is initially trained on DeepSeekMath-Base 7B, employing a learning rate of $2 \times 10^{-5}$. For the GRPO algorithm, we set the policy model’s learning rate to $1 \times 10^{-6}$, and the KL divergence penalty coefficient to $0.04$. For each question, we generate $64$ responses, utilize a maximum sequence length of $1024$ tokens, and adopt a batch size of $1024$ during training. The policy model undergoes only a single update after each exploratory round. Evaluation of DeepSeekMath-RL 7B is conducted on the same benchmarks used by DeepSeekMath-Instruct 7B. In this context, chain-of-thought evaluations on GSM8K and MATH are classified as in-domain tasks for DeepSeekMath-RL 7B, while all remaining benchmarks are considered out-of-domain.

Table \ref{tab:sft_rl_math} presents a comparison of open- and closed-source models employing both chain-of-thought and tool-based reasoning approaches on benchmarks in English and Chinese. The findings highlight: 1) When leveraging chain-of-thought reasoning, DeepSeekMath-RL 7B achieves 88.2\% accuracy on GSM8K and 51.7\% on MATH. This result not only exceeds all other open-source models spanning the 7B to 70B parameter range, but also surpasses most closed-source counterparts. 2) Notably, the only instruction tuning data used for training DeepSeekMath-RL 7B consists of chain-of-thought-format examples from GSM8K and MATH, and it is initialized from DeepSeekMath-Instruct 7B. Although its training dataset is narrowly focused, DeepSeekMath-RL 7B outperforms DeepSeekMath-Instruct 7B on all evaluation measures, thus underscoring the substantial benefits of reinforcement learning.

\section{Discussion}
Here, we present and analyze the outcomes of both our pre-training phase and the reinforcement learning experiments.

\subsection{Lessons Learnt in Pre-Training}

We begin by outlining our observations during pre-training. Except where noted, the training configuration aligns with that described in Section \ref{sec:quality-policy}. Additionally, it should be emphasized that, within this section, references to the DeepSeekMath Corpus pertain to an 89B-token dataset obtained from the second round of data gathering.

\subsubsection{Code Training Benefits Mathematical Reasoning}

An often-cited but unproven hypothesis posits that training on code enhances reasoning skills. In this work, we seek to partially address this claim, focusing on mathematics: code training enhances a model’s capacity for mathematical reasoning, regardless of whether external tools are employed.

To study how code training affects mathematical reasoning, we experimented with the following two-stage training and one-stage training settings:

\textbf{Two-Stage Training}

\begin{itemize}[topsep=0pt]
    \item \textbf{Code Training for 400B Tokens $\rightarrow$ Math Training for 150B Tokens}: The DeepSeek-LLM 1.3B model is initially pre-trained on 400B code tokens, after which it undergoes continued training on an additional 150B math tokens;
    \item \textbf{General Training for 400B Tokens $\rightarrow$ Math Training for 150B Tokens}: To serve as a comparison, we conduct an experiment where the model is first trained on 400B general tokens (sourced from a large-scale general-purpose corpus constructed by DeepSeek-AI), before the same math-focused training. This design aims to assess the effectiveness of code data relative to general data in bolstering mathematical reasoning.
\end{itemize}

\textbf{One-Stage Training}

\begin{itemize}[topsep=0pt]
    \item \textbf{Math Training Across 150B Tokens}: DeepSeek-LLM 1.3B undergoes training using 150B math tokens;
    \item \textbf{Combined Training Using 400B Code Tokens and 150B Math Tokens}: We observe that performing math training after code training negatively impacts coding abilities. To address this, we explore joint training, blending code and math tokens in a single phase, to evaluate whether this strategy can both enhance mathematical reasoning and reduce the risk of catastrophic forgetting.
\end{itemize}

\paragraph{Results}
The downstream performance across various training configurations is presented in Table \ref{tab:code-math} and Table \ref{tab:code-math-general-eval}.

Program-aided mathematical reasoning derives clear advantages from code training in both two-stage and one-stage training frameworks. Referencing Table \ref{tab:code-math}, code training by itself already leads to considerable gains in solving GSM8K and MATH problems with Python under a two-stage approach. Introducing math training in the second phase brings about additional progress. Notably, in the one-stage paradigm, simultaneous integration of code and math tokens alleviates the problem of catastrophic forgetting observed in two-stage setups, while also improving overall coding capabilities (as illustrated in Table \ref{tab:code-math-general-eval}) and enhancing performance in program-aided mathematical reasoning (see Table \ref{tab:code-math}).

Training with code additionally enhances mathematical reasoning capabilities even in the absence of tool usage. Within the two-stage training framework, the preliminary code training phase yields noticeable improvements. Furthermore, this stage facilitates more effective later math training, culminating in superior overall outcomes. In contrast, merging code tokens and math tokens for a single-stage training process hinders mathematical reasoning when tools are not employed. A possible explanation is that the relatively small size of DeepSeek-LLM 1.3B restricts its ability to comprehensively integrate code and mathematical datasets at the same time.

\subsubsection{ArXiv Papers Seem Ineffective in Improving Mathematical Reasoning}

ArXiv papers are commonly included as a component of math pre-training data \citep{minerva,gpt-f,llemma,mathpile}. However, detailed analysis regarding their impact on mathematical reasoning has not been extensively conducted. Perhaps counter-intuitively, according to our experiments, arXiv papers seem ineffective in improving mathematical reasoning. We experiment with models of different sizes, including DeepSeek-LLM 1.3B and DeepSeek-Coder-Base-v1.5 7B \citep{deepseek-coder}, using arXiv corpora that underwent varied processing pipelines:

\begin{itemize}[topsep=0pt]
    \item \textbf{MathPile} \citep{mathpile}: comprises 8.9 billion tokens, constructed through heuristic-based cleaning and filtering, with scientific arXiv papers accounting for over 85\% of its content;
    \item \textbf{ArXiv-RedPajama} \citep{redpajama}: consists of the full collection of arXiv LaTeX documents, from which preambles, comments, macros, and bibliographies have been stripped, amounting to a total of 28.0 billion tokens.
\end{itemize}

In our study, we independently trained DeepSeek-LLM 1.3B on 150B tokens and DeepSeek-Coder-Base-v1.5 7B on 40B tokens using datasets comprised solely of arXiv papers. The results indicate that utilizing arXiv papers as training data does not significantly enhance the models' mathematical reasoning capabilities. Models exposed exclusively to the arXiv corpus did not exhibit meaningful gains and occasionally performed worse on a range of mathematical evaluations with varying levels of difficulty included in this analysis. The assessments encompass quantitative reasoning benchmarks such as GSM8K and MATH (Table \ref{tab:arxiv-cot}), multiple-choice tasks like MMLU-STEM (Table \ref{tab:arxiv-cot}), and problems in formal mathematics exemplified by miniF2F (Table \ref{tab:arxiv_atp}).

However, this conclusion has its limitations and should be taken with a grain of salt. We have not yet studied:

\begin{itemize}[topsep=0pt]
    \item How arXiv tokens influence particular mathematics-oriented tasks that were beyond the scope of this study, including activities like transforming formal theorems or proofs into more informal presentations;
    \item The potential interactions and outcomes when arXiv tokens are integrated with additional datasets;
    \item If scaling up model size would further enhance the advantages gained from incorporating arXiv papers.
\end{itemize}

Thus, further exploration is required, which we leave for future studies.

\subsection{Insights of Reinforcement Learning}
\subsubsection{Towards to a Unified Paradigm} Here, we introduce a comprehensive framework for examining various training strategies, including SFT, RFT, DPO, PPO, and GRPO, and present experiments to investigate the key components of this unified paradigm. In most cases, the gradient with respect to the parameter $\theta$ for a given training approach may be expressed as:

\begin{equation}
    \nabla_{\theta}\mathcal{J}_{\textcolor{red}{\mathcal{A}}}(\theta) = \mathbb{E}[\underbrace{(q,o) \sim \textcolor{red}{\mathcal{D}}}_{Data \ Source}]\left( \frac{1}{|o|} \sum_{t=1}^{|o|}  \underbrace{GC_{{\mathcal{A}}}(q, o, t, \textcolor{red}{\pi_{{rf}}})}_{Gradient \ Coefficient}  \nabla_{\theta}\log \pi_{\theta}(o_t | q, o_{<t})\right).
\label{eq:objective}
\end{equation}

There exist three key components: 1) \textit{Data Source $\mathcal{D}$}, which determines the training data; 2) \textit{Reward Function $\pi_{{rf}}$}, which is the source of the training reward signal; 3) \textit{Algorithm $\mathcal{A}$}: which processes the training data and the reward signal to the gradient coefficient $GC$ that determines the magnitude of the penalty or reinforcement for the data. We analyze several representative methods based on such a unified paradigm:

\begin{itemize}
    \item  \textbf{Supervised Fine-tuning (SFT)}: In SFT, a pretrained model undergoes fine-tuning on a dataset that is specifically chosen by humans for this purpose.
    \item \textbf{Rejection Sampling Fine-tuning (RFT)}: RFT involves additional fine-tuning of the model trained with SFT. It does so using samples generated by the SFT model in response to SFT questions, but only retains samples whose answers are deemed correct.
    \item \textbf{Direct Preference Optimization (DPO)}: DPO builds upon the SFT model, fine-tuning it with enhanced outputs obtained from the SFT model itself. The training utilizes a pair-wise DPO loss function.
    \item \textbf{Online Rejection Sampling Fine-tuning (Online RFT)}: Unlike traditional RFT, Online RFT initializes its policy model with the SFT model, then proceeds to fine-tune this model by leveraging augmented samples that are dynamically produced by the current policy model during training.
    \item \textbf{PPO/GRPO}: For PPO/GRPO, the policy model is initially set as the SFT model and subsequently optimized by reinforcing it with outputs sampled from the real-time, evolving policy model.
\end{itemize}

Table \ref{tab:data_policy} provides an overview of the constituent elements of these approaches. For a comprehensive exposition of the derivation procedures, see Appendix \ref{app:analysis-rl}.

\paragraph{Observation about Data Source} The data sources are classified into two types: online sampling and offline sampling. Online sampling refers to collecting training data generated by the real-time actions of the current training policy model. In contrast, offline sampling involves using training data obtained from the initial SFT model's samples. Methods such as RFT and DPO utilize offline sampling, whereas Online RFT and GRPO employ online sampling.

Figure \ref{fig:rl-analysis} illustrates that Online RFT achieves substantially better performance than RFT across both benchmarks. Notably, while Online RFT and RFT show similar results during the initial phase of training, Online RFT quickly secures a marked lead as training progresses, highlighting the effectiveness of the online approach. This can be explained by the fact that, early in training, the actor and SFT model behave similarly, and the collected samples show minimal divergence. In contrast, during later training stages, the actor’s sampled outputs diverge more prominently, at which point the benefit of dynamic, real-time sampling becomes much more pronounced.

\paragraph{Observation about Gradient Coefficient} In the algorithm, the reward signal influences the gradient coefficient, which is subsequently used to modify the model parameter. Within our experimental setup, the reward function is partitioned into `Rule' and `Model' components. The `Rule' category evaluates response quality through assessment of answer correctness, whereas the `Model' type involves training a dedicated reward model that assigns a score to each response, drawing on data labeled according to rule-based assessments. As demonstrated by Equations \ref{eq:GC-RFT} and \ref{eq:GC-GRPO}, there is a fundamental distinction between Online RFT and GRPO: only GRPO modulates its gradient coefficient as a function of the reward received from the reward model. This mechanism enables the algorithm to apply varying strengths of positive or negative reinforcement depending on the magnitude of the reward. Conversely, Online RFT is devoid of such adaptability; it offers no penalties for incorrect answers and treats all correct responses with an equal degree of reinforcement, irrespective of their individual qualities.

As illustrated in Figure \ref{fig:rl-analysis}, GRPO outperforms online RFT, underscoring the effectiveness of modifying the coefficients for positive and negative gradients. Moreover, GRPO+PS achieves better results than GRPO+OS, demonstrating the advantage of implementing step-aware, fine-grained gradient coefficients. Additionally, we investigate the effect of iterative RL by carrying out two iterations in our experiments. Figure \ref{fig:iter-rl} reveals that iterative RL markedly enhances performance, with the most substantial gains observed after the first iteration.

\subsubsection{Why RL Works?} This study applies reinforcement learning to a selected subset of instruction tuning data and observes that it yields notable gains over the baseline instruction tuning model. To shed light on the reasons underlying the effectiveness of reinforcement learning, we compare the Pass@K and Maj@K accuracies of both Instruct and RL-augmented models across two benchmarks. The results, depicted in Figure \ref{fig:maj-pass}, reveal that while RL improves Maj@K, it does not lead to corresponding gains in Pass@K. This pattern suggests that the benefit of RL mainly arises from strengthening the distributional robustness of the outputs. Put differently, \textbf{the observed advancement appears to stem from increasing the likelihood of correct answers surfacing in the TopK rather than genuine advances in core model competencies.} This phenomenon aligns with observations from \citep{wang2023making}, where a \textbf{misalignment problem} was documented for reasoning tasks handled by SFT models. Prior work further demonstrates that reasoning ability in SFT models can be raised by deploying preference alignment techniques \citep{yuan2023rrhf,song2023preference,wang2023making}.

\subsubsection{How to Achieve More Effective RL?}
\label{sec:effec_rl}
Our results show that RL performs effectively for tasks involving mathematical reasoning. Additionally, we introduce a comprehensive framework that offers a cohesive perspective on several prominent training approaches. In this framework, each method can be categorized as a form of direct or approximate RL strategy. As encapsulated in Equation \ref{eq:objective}, the framework is structured around three principal elements: Data Source, Algorithm, and Reward Function. Below, we outline prospective research avenues related to these components.

\paragraph{Data Source} The foundation of all training approaches lies in the data source. Within the domain of RL, we define the data source as the set of unlabeled questions, paired with corresponding outputs generated by the policy model. In this study, our approach is limited to utilizing the questions from the instruction tuning phase, with outputs sampled through a straightforward nucleus sampling method. We hypothesize that this limitation may contribute to our RL pipeline's enhancement being restricted to the Maj@K metric. Moving forward, we aim to extend our RL pipeline to questions that are outside the training distribution, while integrating \textbf{advanced sampling (decoding) strategies}, such as those employing tree-search algorithms \citep{yao2023tree}. Moreover, \textbf{efficient inference techniques} \citep{xia-etal-2023-speculative,leviathan2023fast,kwon2023efficient,xia2024unlocking}, which are critical for the policy model's exploration effectiveness, are also of considerable importance.

\paragraph{Algorithms} Algorithms utilize both the data and reward signals to determine the gradient coefficient, subsequently adjusting the model parameters. According to Equation \ref{eq:objective}, current approaches tend to completely \textbf{TRUST} the reward function's signal when modifying the conditional probability for a particular token, either raising or lowering it as needed. Nevertheless, it is unrealistic to assume that the reward signal remains consistently accurate, particularly for highly intricate tasks. For instance, the PRM800K datasets \citep{lightman2023let}, which underwent extensive annotation by trained professionals, still exhibit roughly 20\% incorrect labels\footnote{\url{https://github.com/openai/prm800k/issues/12\#issuecomment-1728491852}}. In response, we aim to investigate reinforcement learning algorithms designed to remain effective in the presence of noisy reward signals. We contend that such \textbf{WEAK-TO-STRONG} \citep{burns2023weak} alignment methodologies have the potential to significantly reshape future learning algorithms.

\paragraph{Reward Function} The reward function serves as the principal source of training feedback. Within RL, this function is typically implemented using a neural reward model. We identify three key avenues for advancing reward models: 1) \textbf{Improving the generalization capacity of the reward model.} The reward model should be able to generalize well, managing both out-of-distribution queries and sophisticated output generations; in the absence of this ability, reinforcement learning risks simply reinforcing the status quo of LLM distributions instead of driving meaningful advancement; 2) \textbf{Capturing and expressing the uncertainty within the reward model.} Accounting for model uncertainty might enable improved interaction between limited reward models and weak-to-strong learning strategies; 3) \textbf{Developing process reward models with both high quality and construction efficiency} so that granular feedback can be provided throughout the reasoning steps \citep{lightman2023let,wang2023math}.

\section{Conclusion, Limitation, and Future Work} 
In this work, we introduce DeepSeekMath, a model that achieves superior results among open-source approaches on the competition-level MATH benchmark, attaining a performance comparable to that of proprietary systems. DeepSeekMath~builds upon DeepSeek-Coder-v1.5 7B and is subjected to continual training over 500B tokens, including a substantial segment—120B math-specific tokens—curated from Common Crawl data. Our detailed ablation studies reveal that web content is a highly valuable source for high-quality mathematical data, whereas arXiv does not contribute as much as initially anticipated. We propose Group Relative Policy Optimization (GRPO), an adaptation of Proximal Policy Optimization (PPO), capable of significantly enhancing mathematical reasoning while reducing memory requirements. Empirical findings indicate that GRPO yields tangible benefits, even when DeepSeekMath-Instruct 7B achieves elevated benchmark scores. Additionally, we propose a unified framework to interpret a range of existing strategies, and discuss several promising avenues for more efficient reinforcement learning.

While DeepSeekMath~demonstrates strong performance on quantitative reasoning tasks, its abilities in areas such as geometry and theorem-proving are not as advanced as those of proprietary models. For example, during preliminary assessments, the model struggled with questions pertaining to triangles and ellipses, suggesting possible biases introduced during data selection in both pre-training and fine-tuning phases. Furthermore, due to limitations in model size, DeepSeekMath~underperforms compared to GPT-4 when it comes to few-shot learning; whereas GPT-4 benefits from few-shot input and exhibits notable improvement, DeepSeekMath~maintains nearly the same performance between zero-shot and few-shot scenarios. Looking ahead, we aim to enhance our data engineering workflow to curate a higher-quality pre-training dataset. Additionally, we plan to investigate further avenues (Section \ref{sec:effec_rl}) for optimizing the use of reinforcement learning in LLMs.

\end{document}